% This is the ADASS_template.tex LaTeX file, 24 Oct 2025.
% It is based on the ASP general author template file, but modified to reflect the specific
% requirements of the ADASS proceedings.

% This is the aspauthor.tex LaTeX file
% Copyright 2021, Astronomical Society of the Pacific Conference Series
% Revision:  29 November 2021

% To compile, at the command line positioned at this folder, type:
% pdflatex ADASS_template
% pdflatex ADASS_template
% This will create a file called ADASS_template.pdf

\documentclass[11pt,twoside]{article}

% Do NOT use ANY packages other than asp2023. 
\usepackage{asp2023}

\aspSuppressVolSlug
\resetcounters

% References must all use BibTeX entries in a .bibfile.
% References must be cited in the text using \citet{} or \citep{}.
% Do not use \cite{}.
% See ManuscriptInstructions.pdf for more details
\bibliographystyle{asp2021}

% The ``markboth'' line sets up the running heads for the paper.
% 1 author: "Surname"
% 2 authors: "Surname1 and Surname2"
% 3 authors: "Surname1, Surname2, and Surname3"
% >3 authors: "Surname1 et al."
% Replace ``Short Title'' with the actual paper title, shortened if necessary.
% Use mixed case type for the shortened title
% Ensure shortened title does not cause an overfull hbox LaTeX error
% See ASPmanual2010.pdf 2.1.4  and ManuscriptInstructions.pdf for more details
\markboth{A. Rizzo et al.}{Quantum Annealing for Gamma-Ray Analysis}

\begin{document}

\title{Quantum Annealing for Gamma-Ray Source Detection and Localization with AGILE Data}

% Note the position of the comma between the author name and the 
% affiliation number.
% Authors surnames should come after first names or initials, eg John Smith, or J. Smith.
% Author names should be separated by commas.
% The final author should be preceded by "and".
% Affiliations should not be repeated across multiple \affil commands. If several
% authors share an affiliation this should be in a single \affil which can then
% be referenced for several author names. If only one affiliation, no footnotes are needed.
% See ManuscriptInstructions.pdf and ASP's manual2010.pdf 3.1.4 for more details
\author{Alessandro~Rizzo,$^1$ Nicol\`o~Parmiggiani,$^2$ Andrea~Bulgarelli,$^2$
Gabriele~Panebianco,$^2$ Carlo~Burigana,$^3$ Vincenzo~Fabrizio~Cardone,$^4$
Farida~Farsian,$^1$ Massimo~Meneghetti,$^2$ Giuseppe~Murante,$^5$
Ghulam~Mustafa,$^2$ Giuseppe~Sarracino,$^4$ Roberto~Scaramella,$^4$
Francesco~Schillir\`o,$^1$ Vincenzo~Testa,$^4$ Tiziana~Trombetti,$^3$
Fabio~Roberto~Vitello,$^1$ C.~Pittori,$^{4,6}$ and M.~Tavani$^{7,8}$}

\affil{$^1$INAF -- Osservatorio Astrofisico di Catania, Via Santa Sofia 78, 95123 Catania, Italy; \email{alessandro.rizzo@inaf.it}}
\affil{$^2$INAF -- Osservatorio di Astrofisica e Scienza dello Spazio, Via Piero Gobetti 93/3, 40129 Bologna, Italy}
\affil{$^3$INAF -- Istituto di Radioastronomia, Via Piero Gobetti 93/3, 40129 Bologna, Italy}
\affil{$^4$INAF -- Osservatorio Astronomico di Roma, Via Frascati 33, 00078 Monteporzio Catone, Italy}
\affil{$^5$INAF -- Osservatorio Astrofisico di Torino, Via Osservatorio 20, 10025 Pino Torinese, Torino, Italy}
\affil{$^6$ASI Space Science Data Center (SSDC), Via del Politecnico snc, 00133 Roma, Italy}
\affil{$^7$INAF -- IAPS, Via del Fosso del Cavaliere 100, I-00133 Roma, Italy}
\affil{$^8$Universit\`a degli Studi di Roma Tor Vergata, Via della Ricerca Scientifica 1, I-00133 Roma, Italy}

% This section is for ADS Processing.  There must be one line per author. paperauthor has 9 arguments.
\paperauthor{Alessandro~Rizzo}{alessandro.rizzo@inaf.it}{}{INAF}{Osservatorio Astrofisico di Catania}{Catania}{Catania}{95123}{Italy}
\paperauthor{Nicol\`o~Parmiggiani}{}{}{INAF}{Osservatorio di Astrofisica e Scienza dello Spazio}{Bologna}{Bologna}{40129}{Italy}
\paperauthor{Andrea~Bulgarelli}{}{}{INAF}{Osservatorio di Astrofisica e Scienza dello Spazio}{Bologna}{Bologna}{40129}{Italy}
\paperauthor{Gabriele~Panebianco}{}{}{INAF}{Osservatorio di Astrofisica e Scienza dello Spazio}{Bologna}{Bologna}{40129}{Italy}
\paperauthor{Carlo~Burigana}{}{}{INAF}{Istituto di Radioastronomia}{Bologna}{Bologna}{40129}{Italy}
\paperauthor{Vincenzo~Fabrizio~Cardone}{}{}{INAF}{Osservatorio Astronomico di Roma}{Monteporzio Catone}{Roma}{00078}{Italy}
\paperauthor{Farida~Farsian}{}{}{INAF}{Osservatorio Astrofisico di Catania}{Catania}{Catania}{95123}{Italy}
\paperauthor{Massimo~Meneghetti}{}{}{INAF}{Osservatorio di Astrofisica e Scienza dello Spazio}{Bologna}{Bologna}{40129}{Italy}
\paperauthor{Giuseppe~Murante}{}{}{INAF}{Osservatorio Astrofisico di Torino}{Pino Torinese}{Torino}{10025}{Italy}
\paperauthor{Ghulam~Mustafa}{}{}{INAF}{Osservatorio di Astrofisica e Scienza dello Spazio}{Bologna}{Bologna}{40129}{Italy}
\paperauthor{Giuseppe~Sarracino}{}{}{INAF}{Osservatorio Astronomico di Roma}{Monteporzio Catone}{Roma}{00078}{Italy}
\paperauthor{Roberto~Scaramella}{}{}{INAF}{Osservatorio Astronomico di Roma}{Monteporzio Catone}{Roma}{00078}{Italy}
\paperauthor{Francesco~Schillir\`o}{}{}{INAF}{Osservatorio Astrofisico di Catania}{Catania}{Catania}{95123}{Italy}
\paperauthor{Vincenzo~Testa}{}{}{INAF}{Osservatorio Astronomico di Roma}{Monteporzio Catone}{Roma}{00078}{Italy}
\paperauthor{Tiziana~Trombetti}{}{}{INAF}{Istituto di Radioastronomia}{Bologna}{Bologna}{40129}{Italy}
\paperauthor{Fabio~Roberto~Vitello}{}{}{INAF}{Osservatorio Astrofisico di Catania}{Catania}{Catania}{95123}{Italy}
\paperauthor{C.~Pittori}{}{}{INAF}{Osservatorio Astronomico di Roma}{Monteporzio Catone}{Roma}{00078}{Italy}
\paperauthor{M.~Tavani}{}{}{INAF}{IAPS}{Roma}{Roma}{00133}{Italy}
% There should be one \aindex line (commented out) for each author. These are used to
% build up the author index for the Proceedings. The surname must come first, followed by
% initials. Note the use of ~ before each initial to control spacing.
% The \author entries (see above) have surname last. These \aindex entries have
% surname first.
% The Aindex.py command willl create them for you after you have constructed the \author
% The first entry should be the first author, for bold-facing the author index page-reference

%\aindex{Rizzo,~A.}
%\aindex{Parmiggiani,~N.}
%\aindex{Bulgarelli,~A.}
%\aindex{Panebianco,~G.}
%\aindex{Burigana,~C.}
%\aindex{Cardone,~V.~F.}
%\aindex{Farsian,~F.}
%\aindex{Meneghetti,~M.}
%\aindex{Murante,~G.}
%\aindex{Mustafa,~G.}
%\aindex{Sarracino,~G.}
%\aindex{Scaramella,~R.}
%\aindex{Schillir\`o,~F.}
%\aindex{Testa,~V.}
%\aindex{Trombetti,~T.}
%\aindex{Vitello,~F.~R.}
%\aindex{Pittori,~C.}
%\aindex{Tavani,~M.}


\begin{abstract}
Maximum Likelihood Estimation (MLE) is the standard method to detect and localize gamma-ray sources in photon count maps. This process requires an optimization over the emission model parameters, so it is a natural candidate for quantum annealing, which encodes an optimization task in the ground state of an Ising Hamiltonian. We reformulate the standard likelihood analysis of AGILE-GRID data as a Quadratic Unconstrained Binary Optimization (QUBO) problem, approximating the Poisson likelihood with a least-squares cost. We solve this task with Simulated Quantum Annealing (SQA) and benchmark it on 200 simulated observations at high Galactic latitude against AGILEpy, the official AGILE analysis library. We found that the mean relative residuals of the source flux are $0.070 \pm 0.128$ for the classical pipeline and $0.070 \pm 0.241$ for the QUBO, while those of the isotropic background are $-0.021 \pm 0.074$ and $0.001 \pm 0.023$, respectively. The QUBO selects a location within one grid step of the true one in 92\% of the maps, among 81 candidate positions on a $0.5^\circ$ grid. These results show that the QUBO formulation reaches an accuracy close to that of the classical pipeline, making this a first step towards applying quantum annealing to more complex problems, where it could offer an advantage over classical optimizers.
\end{abstract}

% These lines show examples of subject index entries. At this stage these have to commented
% out, and need to be on separate lines. Eventually, they will be automatically uncommented
% and used to generate entries in the Subject Index at the end of the Proceedings volume.
% Don't leave these in! - replace them with ones relevant to your paper.
%\ssindex{observatories!space-based!AGILE}
%\ssindex{astronomy!gamma-ray}
%\ssindex{techniques!source detection}
%\ssindex{techniques!model fitting}
%\ssindex{algorithm!optimisation}
%\ssindex{software!framework!agilepy}

% These lines show examples of ASCL index entries. At this stage these have to commented
% out, and need to be on separate lines. Eventually, they will be automatically uncommented
% and used to generate entries in the ASCL Index at the end of the Proceedings volume.
% The ascl.py command will scan your paper on possible code names.
% Don't leave these in! - replace them with ones relevant to your paper.




\section{Introduction}
Gamma-ray sources are associated with the most energetic processes in the universe, and their study is the main scientific goal of missions such as AGILE \citep{Tavani_2009}. These cannot be identified directly in the count map, but must be extracted by fitting a physical model to the observed data. The standard approach treats the observed counts as Poisson variables and maximizes the likelihood function introduced for EGRET \citep{1996ApJ...461..396M}, yielding the source fluxes, the levels of the diffuse backgrounds, and the source position, making this an optimization over the model parameters, and therefore suited for quantum computing. Even though we are still in the Noisy Intermediate-Scale Quantum era, in which devices are limited in size and strongly affected by noise \citep{Preskill_2018}, astrophysics and cosmology are an active line of research in this field \citep{rizzo2025qcnn, rizzo2026comparing, bulgarelli2026qcinaf}. Among the quantum paradigms, quantum annealing is designed specifically for optimization \citep{annealing}.
In this work, we reformulate the classical maximum-likelihood analysis of AGILE-GRID data as a Quadratic Unconstrained Binary Optimization (QUBO) problem, solve it with Simulated Quantum Annealing and benchmark it against AGILEpy, the official AGILE analysis library \citep{agilepy}.

\articlefiguretwo{P204_f1.eps}{P204_f2.eps}{fig:maps}{%
AGILE-GRID maps of the simulated field centered at
$(l,b) = (60.0^\circ, -38.18^\circ)$ in the 100\,MeV--10\,GeV band.
Left: exposure map. Right: example of a simulated count map.}

\section{Methodology}
We validate the method on 200 AGILE-GRID observations simulated with the AGILE Science Tools and the AGILEpy analysis framework \citep{agilepy}. We used the standard AGILE-GRID instrument response function, while the exposure map was derived from a real AGILE observation window, processed over a $50^\circ$ field of view with a $0.5^\circ$ pixel size in the 100\,MeV -- 10\,GeV band. The field is centered at $(l, b) = (60.0^\circ, -38.18^\circ)$, where the Galactic diffuse component is negligible and only the isotropic background and a single point source contribute to the maps. For each map, the true parameters are drawn at random within ranges chosen so that the dataset contains moderately bright sources. An example is shown in Figure \ref{fig:maps}.

\subsection{QUBO Formulation}
In the standard analysis, the counts in each pixel $(i, j)$ are treated as independent Poisson variables with expectation:
\begin{equation}
    \theta_{ij} = g_b B_{ij} + c_a S_{ij}(\alpha, \delta),
    \label{eq:model}
\end{equation}
where $B_{ij}$ and $S_{ij}$ are the templates of the isotropic background and of the source flux scaled respectively by their parameters $g_b$ and $c_a$.
The Poisson log-likelihood cannot be cast directly into a QUBO, because it is not quadratic and its parameters are continuous. We therefore replace it with its Gaussian approximation, so that the negative log-likelihood reduces to the least-squares cost:
\begin{equation}
    \chi^{2} = \sum_{ij}\left(g_{b} B_{ij} + c_{a} S_{ij} - n_{ij}\right)^{2} = \lVert Rx - d \rVert^{2} = x^TWx - 2V^Tx + d^Td,
    \label{eq:cost}
\end{equation}
where $d$ is the flattened count map, $x = (g_b, c_a)^T$ and the columns of $R$ are the emission templates, and $W = R^TR$, $V = R^Td$. Each parameter is expanded in binary form over a chosen range, so that the cost depends only on binary variables, whose biases and couplings are programmed on the qubits of the annealer. Each single bias weighs the counts added by a bit against the alignment of the data with its template, so that a bit is switched on only when the data reward exceeds its cost, while the couplings prevent the same counts from being explained by both components. The ground state is represented by the bit configuration whose decoded values minimize $\chi^2$.

The source position is treated in the same way, we replace the continuous coordinates by a regular grid of $M = 81$ candidate locations, a $9 \times 9$ grid with a $0.5^\circ$ step covering a $4^\circ \times 4^\circ$ region around the source, each carrying its own unit-flux template, flux register and binary indicator $z_m$ and a one-hot penalty that enforces that only one node is selected.
The resulting QUBO is solved with the SQA sampler of OpenJij, a path-integral Monte Carlo emulation of quantum annealing \citep{openjij}, and benchmarked against the AGILEpy maximum-likelihood pipeline.


\section{Results}

\articlefigure{P204_f3.eps}{fig:residuals}{%
Distributions of the relative residuals over the 200 simulated maps for AGILEpy (dark) and the QUBO (light).}

The relative residuals of the source flux have a mean and standard deviation of $0.070 \pm 0.128$ for AGILEpy and $0.070 \pm 0.241$ for the QUBO. The isotropic normalization is recovered with relative residuals of $-0.021 \pm 0.074$ classically, and $+0.001 \pm 0.023$ with the quantum approach. These results are compared in Figure \ref{fig:residuals}. The localization performance of the QUBO is close to that of the classical pipeline, 92\% of the sources are placed within one grid step of their true position, with a mean error of $0.26^\circ$, compared with 98\% of AGILEpy using the same grid.

\section{Conclusion and Future Work}
We have shown that the standard maximum-likelihood analysis of AGILE-GRID data can be reformulated as a QUBO problem and solved by annealing with an accuracy close to that of the classical pipeline, recovering flux, background and position with comparable results. As a next step, we plan to validate the method on a larger dataset of simulated maps and to explore the low-significance regime. We will also extend the model to multiple point sources and to an additional free parameter, that is the Galactic diffuse emission. In the longer term, we will introduce spectral models in the QUBO formulation, so that the spectral parameters of the source can also be estimated. Finally, our goal is to run the same formulation on real quantum annealing hardware such as D-Wave Advantage and Advantage2, where we expect the main limitation to be the analog precision of the hardware biases and couplings.





\acknowledgements
This work is supported by the ICSC - Centro Nazionale di Ricerca in HPC, Big Data e Quantum Computing, CN00000013, spoke 10 ``quantum computing'', CUP C53C22000350006. 
The AGILE Mission has been funded and operated by the Italian Space Agency (ASI) with scientific and programmatic participation by the Italian National Institute for Astrophysics (INAF) and the Italian National Institute for Nuclear Physics (INFN), with the data processing done at the ASI/SSDC in Rome up to 2025. Currently, support for the post-operational phase and ongoing archival data analysis is funded by INAF.



\bibliography{P204}  % For BibTex

% if we have space left, we might add a conference photograph here. Leave commented for now.
% \bookpartphoto[width=1.0\textwidth]{foobar.eps}{FooBar Photo (Photo: Any Photographer)}

\end{document}
